% Diagrama del flujo de trabajo humano + agentes. Compilar: pdflatex workflow.tex
\documentclass[tikz,border=4pt]{standalone}
\usepackage[utf8]{inputenc}
\usepackage[T1]{fontenc}
\usepackage[scaled]{helvet}
\renewcommand{\familydefault}{\sfdefault}
\usetikzlibrary{positioning,arrows.meta,fit,backgrounds,calc}
\definecolor{autor}{HTML}{2A78D6}
\definecolor{claude}{HTML}{EB6834}
\definecolor{codex}{HTML}{1BAF7A}
\definecolor{tinta}{HTML}{0B0B0B}
\definecolor{tinta2}{HTML}{52514E}
\definecolor{caja}{HTML}{F4F3F0}

\begin{document}
\begin{tikzpicture}[font=\footnotesize,node distance=6mm and 10mm,
  box/.style={draw=#1,line width=1.2pt,rounded corners=3pt,fill=white,align=center,inner sep=5pt,text=tinta},
  box/.default=tinta2,
  sub/.style={draw=#1,rounded corners=2pt,fill=white,align=center,inner sep=3pt,font=\scriptsize,text=tinta},
  flow/.style={-{Latex[length=2.2mm]},line width=0.8pt,tinta2},
  cross/.style={{Latex[length=2.2mm]}-{Latex[length=2.2mm]},line width=1.1pt,tinta,dashed},
  gate/.style={font=\scriptsize\itshape,text=tinta2}]

% Autor
\node[box=autor,minimum width=11.4cm] (autor) {\textbf{Tesista (humano en el circuito)}\\[1pt]
  \scriptsize escribe los prompts \textperiodcentered{} objeta la física \textperiodcentered{} aprueba cada \textit{commit} \textperiodcentered{} fusiona los PR \textperiodcentered{} corre el cómputo pesado};

% Contexto compartido
\node[box, below=7mm of autor, minimum width=11.4cm, fill=caja] (ctx) {\textbf{Contexto persistente compartido}\\[1pt]
  \scriptsize \texttt{CLAUDE.md} (reglas + mapa + callejones sin salida) \textperiodcentered{} \texttt{AGENTS.md} $\to$ «leé CLAUDE.md» \textperiodcentered{} notas previas en \texttt{claude\_work/}};

% Agentes
\node[box=claude, minimum width=5.6cm, minimum height=2.9cm, anchor=north] (claude) at ($(ctx.south)+(-3.3,-9mm)$) {};
\node[anchor=north,font=\footnotesize\bfseries,text=tinta] at (claude.north) {Claude Code};
\node[sub=claude, below=6mm of claude.north] (plan) {Opus 5: planifica / analiza};
\node[sub=claude, below=2mm of plan] (exec) {Sonnet 5: ejecuta};
\node[sub=claude, below=2mm of exec] (csub) {subagentes paralelos (3--7)};
\node[sub=claude, below=2mm of csub] (cmem) {memoria privada (Codex no la ve)};

\node[box=codex, minimum width=5.6cm, minimum height=2.9cm, anchor=north] (codex) at ($(ctx.south)+(3.3,-9mm)$) {};
\node[anchor=north,font=\footnotesize\bfseries,text=tinta] at (codex.north) {Codex CLI};
\node[sub=codex, below=6mm of codex.north] (astra) {gpt-6-astra («Astra»)};
\node[sub=codex, below=2mm of astra] (xsub) {subagentes por capítulo (3)};
\node[sub=codex, below=2mm of xsub] (guard) {revisor automático de acciones};

\draw[cross] (claude.east) -- node[above,font=\scriptsize,text=tinta,align=center]{revisión\\cruzada} (codex.west);

% Git
\node[box, anchor=north, minimum width=11.4cm, text width=12.4cm] (vcs) at ($(ctx.south |- claude.south)+(0,-8mm)$) {\textbf{Git}: rama por tarea (\textit{worktree}) $\to$ \textit{push} de la rama $\to$ PR $\to$ \textit{merge} a \texttt{main} \textbf{solo por el tesista}};
\node[box, below=5mm of vcs, minimum width=11.4cm, text width=12.4cm] (art) {\textbf{Artefactos}: notas, scripts reproducibles, borradores de capítulos, figuras (\texttt{claude\_work/})};

% Flechas
\draw[flow] (autor) -- (ctx);
\draw[flow] (claude.north |- ctx.south) -- (claude.north);
\draw[flow] (codex.north |- ctx.south) -- (codex.north);
\draw[flow] (claude.south) -- (claude.south |- vcs.north);
\draw[flow] (codex.south) -- (codex.south |- vcs.north);
\draw[flow] (vcs) -- (art);
\coordinate (R) at ($(ctx.east)+(7mm,0)$);
\coordinate (L) at ($(ctx.west)+(-7mm,0)$);
\draw[flow,autor] (autor.east) -- (autor.east -| R) -- node[pos=0.5,right,gate,text=autor,align=left]{«commit\\and push»\\(explícito)} (vcs.east -| R) -- (vcs.east);
\draw[flow,autor] (art.west) -- (art.west -| L) -- node[pos=0.5,left,gate,text=autor,align=right]{revisa,\\corrige,\\reformula\\el pedido} (autor.west -| L) -- (autor.west);
\end{tikzpicture}
\end{document}
